\documentclass[12pt, a4paper, oneside]{ctexbook}
% 全角符号转换为半角符号
\usepackage{newunicodechar}
\newunicodechar{。}{.}
\newunicodechar{，}{,}
\newunicodechar{：}{:}
\newunicodechar{；}{;}
\newunicodechar{！}{!}
\newunicodechar{（}{(}
\newunicodechar{）}{)}
\usepackage{amsmath, amsthm, amssymb, bm, graphicx, enumitem, hyperref, mathrsfs}
% 去掉超链接的红框
\hypersetup{hidelinks}

\title{{\Huge{\textbf{马原期末重点大题汇总}}}}
\author{作者:实空(网名)}
\date{2024-2025学年第一学期}
\linespread{1.5}

\begin{document}

\maketitle

\pagenumbering{roman}
\setcounter{page}{1}

\tableofcontents
\newpage
\setcounter{page}{1}
\pagenumbering{arabic}

\chapter{马克思主义基本原理期末重点大题汇总}

\section{重点大题}

\subsection{不变资本和可变资本（划分的依据、含义、意义）}

\begin{enumerate}
\item 根据资本在剩余价值生产中所起的不同作用，可以将资本划分为不变资本与可变资本。
\item 不变资本是以生产资料形态存在的资本。
\item 可变资本是用来购买劳动力的那部分资本。
\item \textbf{把资本划分为不变资本和可变资本，进一步揭示了剩余价值的源泉。这表明，剩余价值既不是由全部资本创造的，也不是由不变资本创造的，而是由可变资本雇佣的劳动者创造的。雇佣劳动者的剩余劳动是剩余价值的唯一源泉。这种划分也为确定资本家对雇佣劳动者的剥削程度提供了科学依据。}
\end{enumerate}

\subsection{经济全球化及其影响}

\begin{enumerate}
\item \textbf{经济全球化是指在生产不断发展、科技加速进步、社会分工和国际分工不断深化、生产的社会化和国际化程度不断提高的情况下，世界各国、各地区的经济活动越来越超出某一国家和地区的范围而相互联系、相互依存的过程。}
\item \textbf{经济全球化的表现：}
\begin{enumerate}
\item \textbf{生产全球化。}生产全球化是指随着国际分工进一步深化，生产某些高新技术产品不再由某个国家单独完成，而是多个国家协作完成。
\item \textbf{贸易全球化。}贸易全球化是指\textbf{商品和劳务}在全球范围内的自由流动。
\item \textbf{金融全球化。}金融全球化是指世界各国、各地区在\textbf{金融业务、金融政策}等方面\textbf{相互协调、相互渗透、相互竞争}不断加强，使全球金融市场更加开放、金融体系更加融合、金融交易更加自由的过程。
\end{enumerate}
\item \textbf{经济全球化的动因：}
\begin{enumerate}
\item \textbf{科学技术的进步和生产力的发展为经济全球化提供了坚实的物质基础和根本的推动力。}
\item \textbf{跨国公司的发展为经济全球化提供了适宜的企业组织形式。}
\item \textbf{各国经济体制的变革和国际经济组织的发展是经济全球化的体制与组织保障。}
\end{enumerate}
\item \textbf{经济全球化的影响：}
\begin{enumerate}
\item 经济全球化体现了社会化生产的要求，不仅发达国家从中受益，一些发展中国家在参与经济全球化进程中也得到了快速发展。
\begin{enumerate}
\item 经济全球化为发展中国家\textbf{提供先进技术和管理经验}。
\item 经济全球化为发展中国家\textbf{提供更多的就业机会}。
\item 经济全球化\textbf{推动}发展中国家\textbf{国际贸易的发展}。
\item 经济全球化\textbf{促进}发展中国家\textbf{跨国公司的发展}。
\end{enumerate}
\item 经济全球化也是一把“双刃剑”，它在促进经济发展的同时也带来了一些负面影响：
\begin{enumerate}
\item 发达国家与发展中国家在经济全球化进程中的\textbf{地位和收益不平等、不平衡}。
\item 加剧了发展中国家\textbf{资源短缺和环境污染}。
\item 一定程度上增加了\textbf{经济风险}。
\end{enumerate}
\end{enumerate}
\end{enumerate}

\subsection{矛盾的普遍性和特殊性的内容和辩证关系}

\begin{enumerate}
\item 矛盾的普遍性是指矛盾存在于一切事物当中，存在于一切事物发展过程的始终，旧的矛盾解决了，又产生了新的矛盾，事物始终在矛盾中运动。
\item 矛盾的特殊性是指各个具体事物的矛盾、每一个矛盾的各个方面在发展的不同阶段上各有其特点。
\item 矛盾的普遍性和特殊性是辩证统一的关系：
\begin{enumerate}
\item 矛盾的普遍性即矛盾的共性、矛盾的特殊性即矛盾的个性。矛盾的共性是无条件的、绝对的，矛盾的个性是有条件的、相对的。
\item 任何现实存在的事物的矛盾都是共性和个性的有机统一，共性寓于个性之中，没有离开个性的共性，也没有离开共性的个性。
\item 矛盾的共性和个性相统一的关系，既是客观事物固有的辩证法，也是科学的认识方法。
\item 矛盾的普遍性和特殊性的辩证关系原理是马克思主义基本原理同各国实际相结合的哲学基础。
\end{enumerate}
\end{enumerate}

\subsection{唯物辩证法的否定观（否定之否定规律）的内容}

\begin{enumerate}
\item \textbf{否定是事物自我的否定、自我发展，是事物内部矛盾运动的结果。}
\item \textbf{否定是事物发展的环节，}是旧事物向新事物的转变，是从旧质到新质的飞跃。只有经过否定，旧事物才能向新事物转变。
\item \textbf{否定是新旧事物联系的环节，}新事物孕育产生于旧事物，新旧事物是通过否定环节联系起来的。
\item \textbf{辩证否定的实质是“扬弃”，}是新事物对旧事物既批判又继承，既克服其消极因素又保留其积极因素。
\end{enumerate}

\subsection{感性认知和理性认知的辩证关系}

\begin{enumerate}
\item 感性认识是人们在实践基础上，由感觉器官直接感受到的关于\textbf{事物的现象、事物的外部联系、事物的各个方面}的认识。
\item 理性认识是指人们借助抽象思维，在概括整理大量感性材料的基础上，达到关于\textbf{事物的本质、全体、内部联系和事物自身规律性}的认识。
\item 感性认识和理性认识的关系是辩证统一的：
\begin{enumerate}
\item \textbf{理性认识依赖于感性认识。}
\item \textbf{感性认识有待于发展和深化为理性认识。}
\item \textbf{感性认识和理性认识相互渗透、相互包含。}
\item \textbf{感性认识与理性认识的辩证统一关系是在实践的基础上形成的，也需要在实践中发展。}
\end{enumerate}
\end{enumerate}

\subsection{真理与谬误}

\begin{enumerate}
\item 真理是\textbf{标志主观与客观相符合的哲学范畴}，是对客观事物及其规律的\textbf{正确}反映。
\item 谬误是\textbf{同客观事物及其发展规律相违背的认识}，是对客观事物及其发展规律的\textbf{歪曲}反映。
\item 真理与谬误既对立又统一。
\begin{enumerate}
\item \textbf{真理与谬误相互对立。}
\item \textbf{真理与谬误的对立又是相对的。}
\begin{enumerate}
\item \textbf{真理是具体的。}
\item \textbf{真理又是全面的。}
\item \textbf{在批判谬误中发展真理，是谬误向真理转化的一种形式。}
\end{enumerate}
\end{enumerate}
\end{enumerate}

\subsection{新事物是不可战胜的原因}

\begin{enumerate}
\item \textbf{就新事物与环境的关系而言}，新事物之所以新，是因为有新的要素、结构和功能，它适应已经变化了的环境和条件；旧事物之所以旧，是因为它的各种要素和功能已经不适应环境和客观条件的变化，走向灭亡就成为不可避免的。
\item \textbf{就新事物与旧事物的关系而言}，新事物是在旧事物的“母体”中孕育成熟的，它既否定了旧事物中消极腐朽的东西，又保留了旧事物中合理的、适应新条件的因素，并添加了旧事物所不能容纳的新内容。这也正是新事物在本质上优越于旧事物、具有强大生命力的原因所在。
\item \textbf{在社会历史领域}，新事物是社会上先进的、富有创造力的人们创造性活动的产物，它从根本上符合人民群众的利益和要求，能够得到人民群众的拥护，因而必然战胜旧事物。
\end{enumerate}

\subsection{真理的绝对性和相对性的辩证关系}

\begin{enumerate}
\item 真理是标志主观与客观相符合的哲学范畴，是对客观事物及其规律的正确反映。
\item \textbf{真理的绝对性是指真理主客观统一的确定性和发展的无限性。}
\item \textbf{真理的相对性是指人们在一定条件下对客观事物及其本质和发展规律的正确认识总是有限度的、不完善的。}
\item \textbf{真理的绝对性和相对性的辩证统一。}
\begin{enumerate}
\item \textbf{从真理的两重性上看，真理的绝对性和相对性相互依存，任何真理既是绝对的，又是相对的。}
\item \textbf{从真理的发展上看，人类的认识是一个不断深化的过程，永远处在由真理的相对性走向绝对性、接近绝对性的转化和发展过程中。}
\end{enumerate}
\end{enumerate}

\subsection{为什么实践是检验真理的唯一标准}

这是由真理的本性和实践的特点决定的。
\begin{enumerate}
\item \textbf{从真理的本性看，真理是人们对客观事物及其发展规律的正确反映，它的本性在于主观和客观相符合。}检验真理的标准，既不能是主观认识本身，也不能是客观事物。只有那种能够把主观认识与客观事物联系和沟通起来，从而使人们能够把二者加以比较和对照的东西，才能充当检验真理的标准。具有这种特性的东西，只能是作为主客观联系的桥梁、纽带或“交错点”的社会实践。
\item \textbf{从实践的特点看，实践具有直接现实性。}实践的直接现实性是它的客观实在性的具体表现。实践的直接现实性的品格，是实践能够成为检验真理唯一标准的主要根据。
\item \textbf{在实践检验真理的过程中，逻辑证明可以起到重要的补充作用。}但是，逻辑证明只能回答前提与结论的关系是不是符合逻辑问题，而不能回答结论是不是符合客观实际的问题。所以，逻辑证明也不能取代实践作为检验真理的标准。实践，只有实践，才是检验真理的唯一标准。
\end{enumerate}

\subsection{马克思主义的三个基本组成部分}

\textbf{马克思主义哲学}、\textbf{马克思主义政治经济学}和\textbf{科学社会主义}是其三个基本组成部分。

\subsection{马克思主义的三个直接理论来源}

\textbf{德国古典哲学}、\textbf{英国古典政治经济学}、\textbf{英法空想社会主义}为马克思主义的创立提供了直接的理论来源。

\subsection{马克思主义公开问世的标志是什么}

\textbf{1848年2月《共产党宣言》的发表}

\subsection{哲学基本问题的两部分内容}

\begin{enumerate}
\item \textbf{物质和意识何为第一性；}
\item \textbf{存在和思维、物质、意识是否具有同一性。}
\end{enumerate}

\subsection{意识的能动作用主要表现（意识对物质具有反作用）}

\begin{enumerate}
\item \textbf{意识具有目的性和计划性。}
\item \textbf{意识具有创造性。}
\item \textbf{意识具有指导实践改造客观世界的作用。}
\item \textbf{意识具有调控人的行为和生理活动的作用。}
\end{enumerate}

\subsection{事物的普遍联系}

\begin{enumerate}
\item \textbf{联系是指事物内部之间和事物之间相互影响、相互制约、相互作用的关系。}
\item 联系的\textbf{特点}：
\begin{enumerate}
\item 联系具有\textbf{客观性}。
\item 联系具有\textbf{普遍性}。
\item 联系具有\textbf{多样性}。
\item 联系具有\textbf{条件性}。
\end{enumerate}
\end{enumerate}

\subsection{矛盾的同一性和斗争性的含义及其辩证关系}

\begin{enumerate}
\item 矛盾的同一性是指矛盾着的对立面\textbf{相互依存、相互贯通}的性质和趋势。
\item 矛盾的斗争性是指矛盾着的对立面\textbf{相互排斥、相互分离}的性质和趋势。
\item 矛盾的同一性和斗争性相互联结、相辅相成。没有斗争性就没有同一性，没有同一性也没有斗争性；斗争性寓于同一性之中，同一性通过斗争性来体现。矛盾的同一性是有条件的、相对的，矛盾的斗争性是无条件的、绝对的。矛盾的同一性和斗争性相结合，构成了事物的矛盾运动，推动着事物的变化发展。
\end{enumerate}

\subsection{中国特色社会主义的哲学基础是什么}

\textbf{矛盾的普遍性和特殊性辩证关系原理}

\subsection{实践与认识的关系}

\textbf{实践是认识的基础}

\subsection{实践对认识的决定作用}

\begin{enumerate}
\item \textbf{实践是认识的来源。}
\item \textbf{实践是认识发展的动力。}
\item \textbf{实践是认识的目的。}
\item \textbf{实践是检验认识真理性的唯一标准。}
\end{enumerate}

\subsection{真理的检验标准}

人们认识世界的结果是否为客观真理，需要经过检验才能判定，而检验本身需要有一定的标准，这就是真理的检验标准问题。

\subsection{实践标准的确定性和不确定性}

\begin{enumerate}
\item \textbf{实践标准的确定性即绝对性，是指实践作为检验真理标准的唯一性、归根到底性、最终性，离开实践，再也没有其他公正合理的标准。}
\item \textbf{实践标准的不确定性即相对性，是指实践作为检验真理标准的条件性。}
\begin{enumerate}
\item 任何实践都会受到主客观条件的制约，因而都具有不可能完全证实或驳倒一切认识的局限性；
\item 实践是社会的、历史的实践，由于历史条件的种种限制，实践对真理的检验具有相对性、有限性，表现为具体的实践往往只是在总体上证实认识与它所反映的客观事物是否相符合，而不可能绝对地、永恒地、一劳永逸地予以确证。
\end{enumerate}
\end{enumerate}

\subsection{社会意识的相对独立性}

\begin{enumerate}
\item \textbf{社会意识的相对独立性是指，社会意识在从根本上受到社会存在决定的同时，还具有自身特有的发展形式和发展规律。}
\item \textbf{社会意识与社会存在发展具有不完全同步性和不平衡性。}
\item \textbf{社会意识内部各种形式之间存在相互影响且各自具有历史继承性。}
\item \textbf{社会意识对社会存在具有能动的反作用。}
\item \textbf{社会意识的能动作用是通过指导人们的实践活动实现的。}
\end{enumerate}

\subsection{社会形态更替的统一性和多样性}

\begin{enumerate}
\item 社会历史可划分为五种社会形态：\textbf{原始社会、奴隶社会、封建社会、资本主义社会和共产主义社会}。这五种社会形态的依次更替，是社会历史运动的一般过程和一般规律，表现了社会形态更替的统一性。
\item 有些国家在发展中经历了几种社会形态依次更替的典型过程，也有些国家在发展中超越了一个甚至几个社会形态而跨越式地向前发展；有些国家在历史发展的一定阶段上社会形态性质不够典型，甚至多种社会形态特征交叉渗透；有些国家在一定时期由较为落后的社会形态快速跃进为先进的社会形态，而有些国家的社会形态则长期陷于停滞状态；即使是同一种社会形态，在不同国家也会显现出不同特点。所有这些都体现了社会形态更替的多样性。
\end{enumerate}

\subsection{社会形态更替中的必然性和选择性}

\begin{enumerate}
\item 社会形态更替中的必然性：
\begin{enumerate}
\item \textbf{社会形态更替的必然性主要是指社会形态依次更替的过程和规律是客观的，其发展的基本趋势是确定不移的。}
\item \textbf{生产力与生产关系矛盾运动的规律性，从根本上规定了社会形态更替的必然性。}
\end{enumerate}
\item 社会形态更替中的选择性：\textbf{社会形态更替的规律的客观性并不否定人们历史活动的能动性，并不排斥人们在遵循社会发展规律的基础上，对于某种社会形态的历史选择性，人们的历史选择性包含三层意思：}
\begin{enumerate}
\item \textbf{社会发展的必然性造成了一定历史阶段社会发展的基本趋势，为人们的历史选择提供了基础、范围和可能性空间。}
\item \textbf{社会形态更替的过程也是一个主观能动性与客观规律性相统一的过程。}
\item \textbf{人们的历史选择性归根结底是人民群众的选择性。}
\end{enumerate}
\end{enumerate}

\subsection{文化在社会发展中的作用是什么}

\begin{enumerate}
\item 文化是推动社会发展的重要力量，文化对社会发展有着巨大的影响。
\item 文化为社会发展提供\textbf{重要思想指引}。
\item 文化为社会发展提供\textbf{精神动力}。
\item 文化为社会发展提供\textbf{凝聚力量}。
\end{enumerate}

\subsection{为什么人民群众在创造历史中起决定作用}

\begin{enumerate}
\item \textbf{人民群众是社会历史的主体，是历史的创造者。}
\begin{enumerate}
\item 从质上看，人民群众是指一切对社会历史发展起推动作用的人。
\item 从量上看，人民群众是指社会人口中的绝大多数。
\end{enumerate}
\item \textbf{在社会历史发展过程中，人民群众起着决定性的作用。}
\begin{enumerate}
\item 人民群众是\textbf{社会物质财富的创造者}。
\item 人民群众是\textbf{社会精神财富的创造者}。
\item 人民群众是\textbf{社会变革的决定力量}。
\end{enumerate}
\end{enumerate}

\subsection{商品的二因素和生产商品的劳动的二重性}

\begin{enumerate}
\item 商品的二因素：
\begin{enumerate}
\item 使用价值。\textbf{使用价值是指商品能够满足人的某种需要的有用性，反映的是人与自然之间的物质关系，是商品的自然属性}。
\item 价值。\textbf{价值是凝结在商品中的无差别的一般人类劳动，即人的脑力和体力的耗费。价值是商品所特有的社会属性}。
\item 商品的使用价值和价值之间是对立统一的关系。
\end{enumerate}
\item 生产商品的劳动的二重性：
\begin{enumerate}
\item \textbf{具体劳动}。\textbf{具体劳动是指生产一定使用价值的具体形式的劳动。}
\item \textbf{抽象劳动。抽象劳动是指撇开一切具体形式、无差别的一般人类劳动，即人的脑力和体力的耗费。}
\item 具体劳动和抽象劳动也是对立统一的关系。其中劳动的二重性决定了商品的二因素。
\end{enumerate}
\end{enumerate}

\subsection{资本主义生产的绝对规律和作用是什么}

\begin{enumerate}
\item \textbf{生产剩余价值是资本主义生产方式的绝对规律。}
\end{enumerate}

\subsection{剩余价值生产的两种基本方式}

\begin{enumerate}
\item 绝对剩余价值的生产。绝对剩余价值是指在\textbf{必要劳动时间不变}的条件下，由于\textbf{延长工作日的长度或提高劳动强度}而生产的剩余价值。
\item 相对剩余价值的生产。相对剩余价值是指在\textbf{工作日长度不变}的条件下，通过\textbf{缩短必要劳动时间而相对延长剩余劳动时间}所产生的剩余价值。
\end{enumerate}

\subsection{资本主义的基本矛盾}

\textbf{生产社会化和生产资料资本主义私人占有之间的矛盾}，是资本主义的基本矛盾。

\subsection{两个“必然”和两个“决不会”}

\begin{enumerate}
\item 两个“必然”：\textbf{资本主义必然灭亡，社会主义必然胜利}。
\item 两个“决不会”：\textbf{无论哪一个社会形态，在它所能容纳的全部生产力发挥出来以前，是决不会灭亡的}；\textbf{而新的更高的生产关系，在它的物质存在条件在旧社会的胎胞里成熟以前，是决不会出现的}。
\end{enumerate}

\subsection{社会主义发展道路多样性的原因及如何探索社会主义发展道路多样性}

\begin{enumerate}
\item 社会主义发展道路多样性的原因：
\begin{enumerate}
\item \textbf{各个国家的生产力发展状况和社会发展阶段决定了社会主义发展道路具有不同的特点。}
\item \textbf{历史文化传统的差异性是造成不同国家社会主义发展道路多样性的重要条件。}
\item \textbf{时代和实践的不断发展，是造成社会主义发展道路多样性的现实原因。}
\end{enumerate}
\item 探索社会主义（适合本国国情的）发展道路的方法：
\begin{enumerate}
\item 探索社会主义发展道路，必须\textbf{坚持对待马克思主义的科学态度}。
\item 探索社会主义发展道路，必须\textbf{从当时当地的历史条件出发，坚持“走自己的路”}。
\item 探索社会主义发展道路，必须\textbf{充分吸收人类一切优秀文明成果}。
\end{enumerate}
\end{enumerate}

\subsection{共产主义社会的基本特征}

\begin{enumerate}
\item \textbf{物质财富极大丰富，消费资料按需分配。}
\item \textbf{社会关系高度和谐，人们精神境界极大提高。}
\item \textbf{实现每个人自由而全面的发展，人类从必然王国向自由王国飞跃。}
\end{enumerate}

\subsection{社会基本矛盾是历史发展的根本动力，表现有}

\begin{enumerate}
\item \textbf{生产力是社会基本矛盾运动中最基本的动力因素，是人类社会发展和进步的最终决定力量。}
\item \textbf{社会基本矛盾特别是生产力和生产关系的矛盾，决定着社会中其他矛盾的存在和发展。}
\item \textbf{社会基本矛盾具有不同的表现形式和解决方式，并从根本上影响和促进社会形态的变化和发展。}
\end{enumerate}

\subsection{人类社会的基本矛盾和基本规律}

\begin{enumerate}
\item 生产力与生产关系、经济基础与上层建筑之间的矛盾是人类社会的基本矛盾。
\item 生产力与生产关系矛盾运动的规律、经济基础与上层建筑矛盾运动的规律是人类社会的基本规律。
\end{enumerate}

\subsection{产业资本在循环过程中经历的三个不同阶段以及每个阶段的职能}

\begin{enumerate}
\item 购买阶段，执行货币资本的职能。
\item 生产阶段，执行生产资本的职能。
\item 售卖阶段，执行商品资本的职能。
\end{enumerate}

\subsection{产业资本运动的基本前提条件}

产业资本的三种职能形式必须在\textbf{空间上并存，在时间上继起}。

\subsection{垄断并不能消除竞争，反而使竞争变得更加复杂和激烈，原因是什么}

\begin{enumerate}
\item 垄断没有消除产生竞争的经济条件。竞争是商品经济的一般规律。
\item 垄断必须通过竞争来维持。
\item 社会生产是复杂多样的，任何垄断组织都不可能把包罗万象的社会生产全都包下来。
\end{enumerate}

\subsection{量变和质变的辩证关系}

\begin{enumerate}
\item \textbf{量变是质变的必要准备。}
\item \textbf{质变是量变的必然结果，并为新的量变开辟道路。}
\item \textbf{量变和质变是相互渗透的。一方面，在总的量变的过程中有阶段性和局部性的部分质变；另一方面，在质变过程中也有旧质在量上的收缩和新质在量上的扩张。}
\end{enumerate}

\subsection{垄断条件下竞争的特点}

\begin{enumerate}
\item \textbf{在竞争目的上}，自由竞争主要是为获得更多的利润或超额利润，不断扩大资本的积累，而垄断条件下的竞争则是为获取高额垄断利润，并不断巩固和扩大自己的垄断地位和统治权力。
\item \textbf{在竞争手段上}，自由竞争主要运用经济手段，如通过改进技术、提高劳动生产率、降低产品成本等以战胜对手，而垄断条件下的竞争除了采取各种形式的经济手段外，还采取非经济的手段，使竞争变得更加复杂、激烈。
\item \textbf{在竞争范围上}，自由竞争时期，竞争主要是在经济领域，而且主要是在国内市场上进行的，而在垄断时期，国际市场上的竞争越来越激烈，不仅经济领域的竞争多种多样，而且还扩大到经济领域以外。
\end{enumerate}

\subsection{资本主义变化及其实质}

\begin{enumerate}
\item 资本主义变化的主要表现：
\begin{enumerate}
\item \textbf{生产资料所有制的变化。}
\item \textbf{劳资关系和分配关系的变化。}
\item \textbf{社会阶层和阶级结构的变化。}
\item \textbf{经济调节机制和经济危机形态的变化。}
\item \textbf{政治制度的变化。}
\end{enumerate}
\item 资本主义变化的原因：
\begin{enumerate}
\item \textbf{科学技术革命和生产力的发展，是资本主义发生变化的根本推动力量。}
\item \textbf{工人阶级争取自身权力和利益的斗争，是推动资本主义发生变化的重要力量。}
\item \textbf{社会主义制度初步显示的优越性对资本主义产生了重要影响。}
\item \textbf{主张改良主义的政党对资本主义制度的改革，也对资本主义发生变化起到了重要作用。}
\end{enumerate}
\item 资本主义变化的实质：
\begin{enumerate}
\item 第二次世界大战后资本主义发生的变化\textbf{从根本上说是人类社会发展一般规律和资本主义经济规律作用的结果。}
\item 第二次世界大战后资本主义发生的变化\textbf{是在资本主义制度基本框架内的变化，并不意味着资本主义生产关系的根本性质发生了变化。}
\end{enumerate}
\end{enumerate}

\subsection{劳动力成为商品和货币转化为资本的特点与条件分别是什么}

\begin{enumerate}
\item \textbf{劳动力成为商品的基本条件}：
\begin{enumerate}
\item 劳动者在法律上是自由人，能够把自己的劳动力当作自己的商品来支配。
\item 劳动者没有任何生产资料，没有生活资料来源，因而不得不依靠出卖劳动力为生。
\end{enumerate}
\item \textbf{劳动力商品的特点}：
\begin{enumerate}
\item 劳动力的价值和使用价值不同于普通商品。
\item 劳动力的价值，是由\textbf{生产、发展、维持和延续}劳动力所必需的生活必需品的价值决定的。它包括三个部分：
\begin{enumerate}
\item 维持\textbf{劳动者本人}生产所必需的生活资料的价值。
\item 维持\textbf{劳动者家属}的生产所必需的生活资料的价值。
\item 劳动者\textbf{接受教育和训练}所支出的费用。
\end{enumerate}
\item 劳动力的使用价值是价值的源泉，它在消费过程中能够创造新的价值，而且这个新的价值比劳动力本身的价值更大。
\end{enumerate}
\item \textbf{货币转化为资本的条件}：货币所有者购买到劳动力以后，在消费过程中，不仅能够收回他在购买这种特殊商品时支付的价值，还能得到一个增殖的价值即剩余价值。而一旦货币购买的劳动力带来剩余价值，货币也就变成了资本。
\end{enumerate}

\subsection{生产力与生产关系的辩证（相互）关系}

\begin{enumerate}
\item 生产力决定生产关系。
\begin{enumerate}
\item 生产力\textbf{状况}决定生产关系的\textbf{性质}。
\item 生产力的\textbf{发展}决定生产关系的\textbf{变化}。
\end{enumerate}
\item 生产关系对生产力具有能动的反作用。主要表现在：
\begin{enumerate}
\item 当生产关系\textbf{适合}生产力发展的客观要求时，对生产力的发展起\textbf{推动}作用。
\item 当生产关系\textbf{不适合}生产力发展的客观要求时，就会\textbf{阻碍}生产力的发展。
\end{enumerate}
\end{enumerate}

\subsection{交往及其作用}

\begin{enumerate}
\item \textbf{交往的定义：交往是唯物史观的重要范畴，指在一定历史条件下的现实的个人、群体、阶级、民族、国家之间的物质和精神上相互往来、相互作用、彼此联系的活动。}
\item \textbf{根据交往内容和交往方式不同，人们把交往划分为各种不同的类型，总体来说，可以将交往划分为物质交往和精神交往。}
\begin{enumerate}
\item \textbf{物质交往是指人们在物质生产实践中发生的交往，物质产品是其交往内容。}
\item \textbf{精神交往是指一定历史条件下，人们在涉及思想、意识、观念、情感和情绪等精神性的领域中进行的交往。}
\end{enumerate}
\item \textbf{交往的作用}（交往是人类实践活动的重要组成部分，对社会生活有着重要的影响）：
\begin{enumerate}
\item \textbf{促进生产力的发展。}
\item \textbf{促进社会关系的进步。}
\item \textbf{促进文化的发展与传播。}
\end{enumerate}
\end{enumerate}

\subsection{守正创新的内涵和意义}

\begin{enumerate}
\item 内涵：
\begin{enumerate}
\item 要坚持守正不动摇。所谓守正，就是\textbf{坚持实事求是，坚持真理性认识，坚持正确政治方向}。
\item 要坚持创新不停步。所谓创新，就是\textbf{坚持解放思想，破除与客观事物进程不相符合的旧观念、旧模式、旧做法，发现和运用事物的新联系、新属性、新规律，更有效地认识世界和改造世界}。
\item 守正创新就是坚持守正与创新的辩证统一。
\end{enumerate}
\item 意义：
\begin{enumerate}
\item 守正创新深刻揭示了\textbf{坚持真理与发展真理、坚持马克思主义与发展马克思主义的辩证关系，为推进党和人民事业提供了科学的立场观点方法}。
\item 守正创新深刻揭示了\textbf{“变”与“不变”、继承与发展的辩证统一}。
\item 守正创新彰显着\textbf{党的理论创新和实践创新的时代要求}。
\end{enumerate}
\end{enumerate}

\subsection{社会存在和社会意识的辩证关系是什么}

\begin{enumerate}
\item 含义：
\begin{enumerate}
\item 社会存在是指社会物质生活条件，是社会生活的物质方面，主要包括自然地理环境、人口因素和物质生产方式。
\item 社会意识是社会存在的反映，是社会生活的精神方面。
\end{enumerate}
\item 辩证关系：
\begin{enumerate}
\item 社会存在和社会意识是辩证统一的。社会意识是社会存在的反映，并反作用于社会存在。
\item 社会存在是社会意识内容的客观来源，社会意识是社会物质生活过程及其条件的主观反映。
\item 社会存在决定社会意识，社会意识以理论、观念、心理等形式反映社会存在。这是社会意识对社会存在的依赖性。但社会意识并非消极被动地受制于社会存在，它既依赖于社会存在，又有其相对独立性。
\end{enumerate}
\end{enumerate}

\subsection{社会发展中科技的作用是什么}

\begin{enumerate}
\item 科学技术作为先进生产力的重要标志，是推动社会文明进步的重要力量。
\item 科技革命是推动经济和社会发展的强大杠杆。
\begin{enumerate}
\item 对生产方式产生了深刻影响：
\begin{enumerate}
\item \textbf{改变了社会生产的构成要素。}
\item \textbf{改变了人们的劳动形式。}
\item \textbf{改变了社会经济结构，特别是导致产业结构发生变革。}
\end{enumerate}
\item 对生活方式产生了巨大影响。现代科技革命把人们带入了信息时代，要求人们不断更新和充实知识，以适应时代发展的需要。
\item 促进了思维方式的变革。在现代科技革命条件下，人们获得了新的知识理论结构，能够运用新的理论工具和现代化技术手段去研究一系列新现象、新领域、新课题。
\end{enumerate}
\end{enumerate}

\subsection{无产阶级政党的群众路线是什么}

\begin{enumerate}
\item \textbf{一切为了群众}，这是\textbf{群众路线的基本出发点和最终归宿}，这是\textbf{党的根本宗旨}。
\item \textbf{一切依靠群众}，这是\textbf{群众路线的根本要求}。
\item \textbf{从群众中来}，\textbf{到群众中去}，这是\textbf{无产阶级政党的领导方法}，也是\textbf{群众路线的基本工作方法}。
\end{enumerate}

\section{经典论断与对应原理}

\begin{enumerate}
\item 人在“劳动过程结束时得到的结果，在这个过程开始时就已经在劳动者的表象中存在着，即已经观念地存在着”——\textbf{意识具有目的性和计划性}
\item “世界不会满足人，人决心以自己的行动来改变世界”——\textbf{意识具有指导实践改造客观世界的作用}
\item “当我们通过思维来思考自然界或人类历史或我们自己的精神活动的时候，首先呈现在我们眼前的，是一幅由种种联系和相互作用无穷无尽地交织起来的画面”——\textbf{事物的普遍联系}
\item “社会一旦有技术上的需要，这种需要就会比十所大学更能把科学推向前进”——\textbf{实践是认识发展的动力}
\item “我们党现阶段提出和实施的理论和路线方针政策，之所以正确，就是因为它们都是以我国现时代的社会存在为基础的。”——\textbf{实践是认识的来源}
\item “判定认识或理论之是否真理，不是依主观上觉得如何而定，而是依客观上社会实践的结果如何而定。真理的标准只能是社会的实践。”——\textbf{实践是检验认识真理性的唯一标准}
\item “我们实践证明：感受到了的东西，我们不能立刻理解它，只有理解了的东西才更深刻地感受它”——\textbf{感性认识和理性认识的辩证关系}
\item “对自然界的一切真实的认识，都是对永恒的东西、对无限的东西的认识，因而本质上是绝对的”——\textbf{真理的绝对性}
\item “人不能完全地把握、反映、描绘整个自然界、它的‘直接的总体’，人只能通过创立抽象、概念、规律、科学的世界图景等等永远地接近于这一点。”——\textbf{真理的相对性}
\item “任何真理，如果把它说得‘过火’……加以夸大，把它运用到实际适用的范围之外，便可以弄到荒谬绝伦的地步，而且在这种情形下，甚至必然会变成荒谬绝伦的东西。”——\textbf{真理与谬误的辩证关系}
\item “在这里不要忘记：实践标准实质上绝不能完全地证实或驳倒人类的任何表象。这个标准也是这样的‘不确定’，以便不让人的知识变成‘绝对’，同时它又是这样的确定，以便同唯心主义和不可知论的一切变种进行无情的斗争。”——\textbf{实践标准的确定性与不确定性}
\item “手推磨产生的是封建主的社会，蒸汽磨产生的是工业资本家的社会。”——\textbf{生产力决定生产关系}
\item “一个人的发展取决于和他直接或间接进行交往的其他一切人的发展”——\textbf{交往促进人的全面发展}
\item “世界历史发展的一般规律，不仅丝毫不排斥个别发展阶段在发展的形式或顺序上表现出特殊性，反而是以此为前提的。”——\textbf{社会形态更替的统一性与多样性}
\item “人民，只有人民，才是创造世界历史的动力”——\textbf{人民群众是历史的创造者}
\item “资产阶级的灭亡和无产阶级的胜利是同样不可避免的。”——\textbf{资本主义必然灭亡，社会主义必然胜利}
\item “无论哪一个社会形态，在它所能容纳的全部生产力发挥出来以前，是决不会灭亡的；而新的更高的生产关系，在它的物质存在条件在旧社会的胎胞里成熟以前，是决不会出现的。”——\textbf{两个“决不会”}
\end{enumerate}




\end{document}